\chapter*{Abstract}
\thispagestyle{fancy}
\addcontentsline{toc}{chapter}{Abstract}
\begingroup
\setlength{\emergencystretch}{2em}

This report analyzes normal-incidence electromagnetic reflection from a finite,
lossless dielectric slab placed between two air regions. The slab thickness is fixed
at one free-space wavelength at 1~GHz, approximately 299.79~mm, while the relative
permittivity is swept from 1 to 10 in increments of 0.25. The preserved CST
frequency-domain results are interpreted using the equivalent transmission-line model
of an air--dielectric--air section.

For an ideal slab of thickness $d=\lambda_0$, the half-wavelength matching condition
predicts reflection minima at
$\varepsilon_r=(m/2)^2$. Within the simulated range, the corresponding values are
$\varepsilon_r=1$, 2.25, 4, 6.25, and 9. The CST sweep exhibits minima at the same
permittivity values. Preserved linear-magnitude values include
$|S_{11}|=1.94892\times10^{-4}$ at $\varepsilon_r=1$,
$3.21858\times10^{-3}$ at 2.25,
$1.04206\times10^{-2}$ at 4, and
$1.28988\times10^{-2}$ at 6.25. These correspond to approximately
$-74.20$, $-49.85$, $-39.64$, and $-37.79$~dB, respectively.

Electric- and magnetic-field plots at $\varepsilon_r=4$ illustrate the
TEM-like propagation environment and the matched electrical-length condition.
The agreement in the locations of the reflection minima provides the principal
validation result. Exact null depth is not treated as the primary metric because the
CST model has finite transverse dimensions, numerical ports, finite mesh resolution,
approximate side-boundary representation, and an unresolved higher-order port-mode
warning.

\par\endgroup
